\documentclass[11pt,a4paper]{article}

\usepackage[a4paper,margin=24mm,headheight=15pt]{geometry}
\usepackage{amsmath,amssymb}
\usepackage{booktabs,longtable,tabularx,array,multirow}
\usepackage{graphicx}
\usepackage{float}
\usepackage{caption}
\usepackage{subcaption}
\usepackage{siunitx}
\usepackage{xcolor}
\usepackage{enumitem}
\usepackage{listings}
\usepackage{seqsplit}
\usepackage[numbers,sort&compress]{natbib}
\usepackage{hyperref}
\usepackage[nameinlink,noabbrev]{cleveref}
\usepackage{fancyhdr}

\definecolor{linkblue}{RGB}{25,74,133}
\definecolor{codegray}{RGB}{246,247,249}
\hypersetup{
  colorlinks=true,
  linkcolor=linkblue,
  citecolor=linkblue,
  urlcolor=linkblue,
  pdftitle={Interpretable Detection of AI-Generated and Human Music: Spectral, Dynamic, and Rhythmic Evidence},
  pdfauthor={Experiment archive}
}
\urlstyle{tt}
\captionsetup{font=small,labelfont=bf}
\setlist{nosep,leftmargin=2em}
\sisetup{detect-all,group-separator={,},group-minimum-digits=4}
\setlength{\emergencystretch}{3em}
\setlength{\parindent}{2em}
\setlength{\parskip}{0.25em}
\renewcommand{\arraystretch}{1.18}

\pagestyle{fancy}
\fancyhf{}
\fancyhead[L]{Interpretable Detection of AI Music}
\fancyhead[R]{Updated: 2026-09-04}
\fancyfoot[C]{\thepage}

\newcommand{\coderef}[1]{\hyperlink{code:#1}{\texttt{#1}}}
\newcommand{\outref}[1]{\hyperlink{output:#1}{\texttt{#1}}}
\newcommand{\analysisref}[1]{\hyperlink{analysis:#1}{\texttt{#1}}}
\newcommand{\codeanchor}[1]{\hypertarget{code:#1}{\texttt{#1}}}
\newcommand{\outanchor}[1]{\hypertarget{output:#1}{\texttt{#1}}}
\newcommand{\analysisanchor}[1]{\hypertarget{analysis:#1}{\texttt{#1}}}
\newcommand{\filepath}[1]{\path{#1}}
\newcommand{\hashvalue}[1]{{\ttfamily\seqsplit{#1}}}
\newcommand{\MFFigureRoot}{../../music_flamingo_midi_rhythm_probe_20260903}
\newcommand{\MFHash}[1]{{\ttfamily\seqsplit{#1}}}

\lstset{
  basicstyle=\ttfamily\footnotesize,
  backgroundcolor=\color{codegray},
  frame=single,
  rulecolor=\color{black!15},
  breaklines=true,
  columns=fullflexible,
  keepspaces=true,
  showstringspaces=false,
  upquote=true
}

\title{Interpretable Detection of AI-Generated and Human Music\\
\large Spectral, Dynamic, and Rhythmic Evidence with Reproducible Artifact References}
\author{Undergraduate Thesis Experiment Archive}
\date{Spectral experiment frozen: September 1, 2026\\External dynamics/rhythm/structure validation added: September 3, 2026\\External 500-per-generator test set frozen: September 4, 2026}

\begin{document}

\maketitle

\begin{abstract}
This report extends an existing Human-versus-Suno spectral study to HeartMuLa and ACE-Step 1.5 and adds separately frozen validation of dynamics, rhythm, and structure features. The spectral workflow covers a 500-item conditional manifest, RTX 5090 inference, unified audio standardization, Demucs vocal separation, separator-response correction, and within-domain, cross-generator, prompt-paired, and duration-sensitivity analyses. Vocal stems of all three AI generators deviate from Human music in the 10--16 kHz region, while the 16--20 kHz direction depends on the generator. The highest within-domain locked-test ROC-AUC values are 0.922 for Suno, 1.000 for HeartMuLa, and 0.988 for ACE-Step, but cross-generator transfer remains asymmetric. A MIDI-grounded Groove MIDI Dataset probe rejects Music Flamingo as a quantitative dynamics/rhythm extractor. A subsequent independent fidelity benchmark shows that base beat/downbeat localization and event-level piano dynamics are measurable, but its detector-facing summaries fail external fidelity thresholds. A separately frozen controlled-change benchmark then tests the operational end-to-end question requested for all three families. Non-vocal dynamics passes its external sensitivity gate and increases pooled spectral AUC from 0.8754 to 0.8919, with positive increments for every held-out generator. Rhythm and section/bar-length family gates fail; their AI/Human ablations remain exploratory even when their union improves cohort AUC. Finally, a prompt-paired external test set is frozen with 500 tracks from each of ten non-Suno, non-HeartMuLa, non-ACE-Step generators and 500 shared Human controls; detector evaluation on this set is intentionally pending. Stable file references are supplied for every method, log, output, result table, and figure.
\end{abstract}

\noindent\fcolorbox{black!25}{codegray}{%
\parbox{0.96\linewidth}{\textbf{Scope of the conclusion.} This study establishes separable spectral evidence and a modest incremental non-vocal dynamics signal within the frozen experimental domains; it does not establish a universal AI-music detector. The Human FMA controls are not track-matched to the Muse-conditioned generated songs. Genre, vocal activity, language, mastering, and source provenance may contribute to discrimination. The GMD and GTZAN/MAESTRO/SALAMI experiments test measurement capability or change sensitivity; the subsequent cohort ablation measures association, not causal authorship.}}

\tableofcontents
\listoffigures
\listoftables
\clearpage

\section{Research Questions and Experimental Scope}
\label{sec:scope}

The study begins with a subjective but operational listening observation: vocals, cymbals, sibilants, and other high-frequency transients often give AI-generated music a different sense of ``air,'' grain, bandwidth edge, or over-smoothing from human recordings. This observation is separated into four research questions:

\begin{enumerate}
  \item Can this perceptual cue be captured by interpretable measurements of spectral tilt, band-energy ratios, spectral texture, modulation, and pitch periodicity?
  \item After reducing arrangement-level interference and analyzing only Demucs-separated vocals, do Human--AI differences remain?
  \item Do the differences disappear after applying an externally frozen estimate of the average Demucs frequency-response bias?
  \item When transferred across Suno, HeartMuLa, and ACE-Step, do the measurements describe a general AI cue or a fingerprint of a particular generator or codec?
  \item After external controlled-change validation, do non-vocal dynamics, rhythm change, or section/bar-length variety add locked and generator-held-out information beyond the vocal spectral baseline?
\end{enumerate}

HeartMuLa is an open music-foundation-model family containing HeartCodec and autoregressive song-generation models \citep{heartmula2026,heartlib2026}. ACE-Step 1.5 combines language-model planning with diffusion-based generation \citep{acestep2026,acestepcode2026}. This extension generates 500 items from each system and compares them with the existing 500 Suno and 500 Human items. The frozen experimental protocol and research notes are indexed as \outref{O10} and \outref{O11}.

\section{Experimental Design}
\label{sec:design}

\subsection{Data Domains and Frozen Splits}

The Human controls come from the FMA domain \citep{fma2017}. The AI domains are Suno, HeartMuLa, and ACE-Step. HeartMuLa and ACE-Step use text metadata from the open Muse dataset \citep{muse2026,musedataset2026}. The accompanying Muse Suno V5 audio was \textbf{not downloaded, decoded, or analyzed}; consequently, no Muse waveform enters this experiment.

The frozen prompt manifest was constructed by \coderef{C01}. Its SHA-256 digest is:

\begin{center}
\small\hashvalue{bb0941b91d41d10a978cb9e4b077399e89fea92526ee469e0e6743b09efb917b}
\end{center}

The selection procedure retained 250 Chinese and 250 English examples, required \texttt{style\_sim} $\geq0.4$, required the source metadata to cover at least 30 s, and included only lyric sections within the first 30 s. ACE captions were limited to 512 characters and lyrics to 4,096 characters. Eligible records were ordered by salted SHA-256 so that the selection was independent of source-JSONL ordering. For each language, 200 records were assigned to development and 50 to locked test, yielding a 400/100 split. Every source ID was also mapped to a unique integer seed. The manifest, selected source rows, and metadata are indexed as \outref{O01}.

\begin{table}[H]
\centering
\caption{Frozen cohorts and their roles. Each Human-versus-generator comparison contains 500 Human items. The pooled locked evaluation retains only one copy of each Human item.}
\label{tab:cohort}
\begin{tabular}{lrrrrl}
\toprule
Domain & Total & Development & Locked test & Language control & Role \\
\midrule
Human FMA & 500 & 400 & 100 & Unmatched & Control domain \\
Suno & 500 & 400 & 100 & Existing frozen set & Historical baseline \\
HeartMuLa & 500 & 400 & 100 & 250 ZH/250 EN & New generator \\
ACE-Step & 500 & 400 & 100 & 250 ZH/250 EN & New generator \\
\bottomrule
\end{tabular}
\end{table}

\subsection{Frozen Model and Dataset Versions}

To prevent floating code or dataset revisions from undermining reproducibility, code commits, weight revisions, dependency freezes, and file-level hashes were recorded before inference. See \outref{O06}, \outref{O07}, and \outref{O12}.

\begin{table}[H]
\centering
\scriptsize
\caption{Frozen generator, code, dataset, and weight versions.}
\label{tab:model-versions}
\begin{tabularx}{\linewidth}{lX}
\toprule
Component & Frozen version \\
\midrule
HeartMuLa code & \filepath{HeartMuLa/heartlib@3783bdb8441f2c298b1e64c8651173aac200361c} \\
HeartMuLa 3B weights & revision \filepath{41f6fc68490e11dc43fdabaa6b5767946408c903} \\
HeartCodec & revision \filepath{f889dab0532cfa4bf459f2a3367eb6d346b8eeda} \\
Heart tokenizer/configuration & revision \filepath{9906b2bcd4598772a32cad4aec0760170fe0d177} \\
ACE-Step code & \filepath{ace-step/ACE-Step-1.5@ca1e85fe9430179831e6bc6be790c332190a3866} \\
ACE-Step weights & revision \filepath{19671f406d603126926c1b7e2adc169acbcade22} \\
Muse metadata & revision \filepath{b1bf3bf906daab3a896e14f6dea58cc295848452} \\
Demucs correction & \filepath{demucs_frequency_bias.npz}, SHA-256 \hashvalue{bcacfecac5ce69927dfed2a51ec21cc346f613a7874c519e419d22d35a97de5e} \\
\bottomrule
\end{tabularx}
\end{table}

\subsection{Inference Parameters and Compute Constraints}

Model inference and Demucs separation were performed only on the RTX 5090 GPUs of host \texttt{5090-4}; UVR was not run locally. Before allocating the host, GPU memory, utilization, compute processes, host load, and disk headroom were checked independently on 5090-2, 5090-4, and 5090-5. At that time, all eight GPUs on 5090-4 showed approximately 2 MiB allocated memory and 0\% utilization, so that host was selected as the isolated experimental machine. The environment and selection evidence are preserved in \outref{O06}.

HeartMuLa was invoked through \coderef{C03} with a target duration of 30,000 ms, top-k 50, temperature 1.0, and CFG 1.5. MuLa used bf16, HeartCodec used fp32, and the model and codec remained resident on separate GPUs. ACE-Step was invoked through \coderef{C02} using \texttt{acestep-v15-turbo}, 30 s, eight diffusion steps, and a fixed seed. Thinking/chain-of-thought rewriting and built-in peak normalization were disabled, while the release implementation's DCW setting \texttt{double/0.05/0.02/Haar} was preserved.

\subsection{Unified Audio Protocol}

Raw outputs from both new generators were standardized by \coderef{C05} to:

\begin{itemize}
  \item exactly 30 s;
  \item 44.1 kHz sampling rate;
  \item stereo;
  \item 1,323,000 frames per channel;
  \item PCM-16 FLAC; and
  \item no loudness normalization, equalization, or full-track peak normalization.
\end{itemize}

Standardization consisted only of SoundFile decoding, SciPy \texttt{resample\_poly} polyphase resampling, channel normalization, trimming or zero-padding, and PCM encoding. The few samples that exceeded the valid amplitude range after resampling were clipped individually and counted. \coderef{C06} verified sampling rate, duration, channels, finite values, and file counts. When an atomic rename had succeeded but a controlled interruption prevented a state append, \coderef{C04} reconstructed only the missing state entry and did not rewrite the audio.

\subsection{Demucs Vocal Separation}

Separation used \texttt{htdemucs} from Hybrid Transformer Demucs \citep{demucs2021}, a hybrid waveform-and-spectrogram source-separation model. MUSDB18 is a standard associated separation benchmark \citep{musdb18}. The batch wrapper is \coderef{C07}, with the following frozen parameters:

\begin{lstlisting}[language=bash,caption={Frozen Demucs vocal-only parameters.}]
htdemucs --two-stems vocals --float32 \
  --shifts 0 --overlap 0.25 --segment 7 --jobs 1
\end{lstlisting}

The runtime used Demucs 4.0.1 and PyTorch/Torchaudio 2.7.1+cu128. The read-only environment freeze is in \outref{O06}. After a smoke test, formal separation ran only on GPU 3. Each input produced \texttt{vocals.wav} and \texttt{no\_vocals.wav}; both were 44.1 kHz, stereo, 30 s, float32 WAV.

\section{Spectral Measurements and Statistical Methods}
\label{sec:method}

\subsection{STFT and Spectral Shape}

All metrics were computed on mono vocal stems. Stereo stems were averaged and DC removed. Let $x[m]$ denote the waveform, with a Hann window of length $N=4096$ and hop $H=1024$:
\begin{equation}
X(n,k)=\sum_{m=0}^{N-1}x[nH+m]w[m]\exp\!\left(-j\frac{2\pi km}{N}\right).
\label{eq:stft}
\end{equation}
Power was defined as $P(n,k)=|X(n,k)|^2+\epsilon$. The time-frame rate for a 30 s recording was approximately $44100/1024\approx43.1$ Hz. Average-difference heatmaps used 128 logarithmic bands from 300 Hz to 20 kHz and were generated by \coderef{C12}.

The absolute representation preserved overall level. The frame-normalized representation subtracted the mean over the 300 Hz--20 kHz bands from each frame:
\begin{equation}
S_{\mathrm{shape}}(n,b)=S(n,b)-\frac{1}{|\mathcal{B}|}\sum_{b'\in\mathcal{B}}S(n,b').
\label{eq:shape}
\end{equation}
This operation attenuated global level differences and emphasized the relative distribution of energy along the frequency axis. The population-average difference between AI and Human was defined as
\begin{equation}
\Delta(n,b)=\frac{1}{N_{\mathrm{AI}}}\sum_i S_i^{\mathrm{AI}}(n,b)
-\frac{1}{N_{\mathrm{H}}}\sum_j S_j^{\mathrm{H}}(n,b).
\label{eq:average-diff}
\end{equation}
Songs were not waveform-aligned. Therefore, the horizontal axis describes a shared position in the 30 s structural window and must not be interpreted as sample-wise cancellation.

\subsection{Frozen Demucs Frequency-Response Correction}

Let $b(k)$ dB denote the average Demucs frequency bias estimated on external MUSDB oracle vocals. It was frozen before the current AI/Human labels were examined and applied identically to every vocal stem:
\begin{equation}
X_{\mathrm{corr}}(n,k)=X(n,k)10^{-b(k)/20}.
\label{eq:bias-correction}
\end{equation}
Phase was preserved, followed by ISTFT overlap-add and DC removal. The implementation in \coderef{C11} first performed an STFT/ISTFT identity test with zero bias, requiring relative error no greater than $10^{-8}$. This correction removes only the shared average separator response; it cannot remove class-dependent separation error.

\subsection{Interpretable Scalar Metrics}

Metric computation is centralized in \coderef{C09} and restricted to active frames. The activity threshold for each song was the greater of the 30th percentile of vocal-frame RMS and 35 dB below that song's maximum RMS. If fewer than 16 frames remained, the 16 highest-RMS frames were retained. The metrics were organized into four groups.

\begin{table}[H]
\centering
\small
\caption{Principal vocal spectral features and their perceptual or signal interpretation.}
\label{tab:metrics}
\begin{tabularx}{\linewidth}{p{34mm}p{49mm}X}
\toprule
Group & Representative metrics & Interpretation \\
\midrule
Brightness and tilt & \texttt{hf\_tilt\_5\_16k}, \texttt{hf\_ratio\_5\_16}, \texttt{air\_ratio\_12\_20}, \texttt{sibilance\_ratio} & High-frequency slope, energy relative to the vocal body, air-band energy, and the 5--10 kHz sibilance ratio \\
High-frequency texture & \texttt{hf\_flatness}, \texttt{hf\_entropy}, \texttt{hf\_crest}, \texttt{fakeprint\_periodicity} & Spectral flatness, entropy, peak-to-average contrast, and fixed-frequency or comb-like texture \\
High-frequency dynamics & \texttt{hf\_flux}, \texttt{hf\_mod\_4\_12}, \texttt{sibilance\_burst\_rate} & Inter-frame change, 4--12 Hz high-frequency envelope modulation, and sibilance-burst density \\
Periodicity and pitch & \texttt{CPP}, \texttt{HNR}, \texttt{pitch confidence}, \texttt{vibrato depth} & Glottal periodicity, harmonic-to-noise ratio, pitch stability, and 4--8 Hz vibrato structure \\
\bottomrule
\end{tabularx}
\end{table}

The high-frequency ratio was $10\log_{10}(E_{5-16\mathrm{k}}/E_{1-5\mathrm{k}})$, and the air ratio compared 12--20 kHz with 1--5 kHz. Spectral tilt was the linear slope of mean-band dB against $\log_2 f$. The sibilance-burst threshold was the within-song median sibilance ratio plus 1.5 times its median absolute deviation; burst count was divided by active duration.

\subsection{Frequency-Band Vectors and Classification}

In addition to the 22 scalar metrics, \coderef{C10} constructed a vector for each frequency representation. Each range was divided into 48 geometrically spaced bands. The active-frame mean and standard deviation in dB were recorded for every band, after which the mean components were centered. The ablations were full 0.3--20 kHz, only 5--10 kHz, without 5--10 kHz, only 5--16 kHz, and low-mid 0.3--5 kHz.

The classifier was a standardized ridge linear score with $\alpha=10$:
\begin{equation}
\hat{\mathbf{w}}=(\mathbf{Z}^{\mathsf T}\mathbf{Z}+\alpha\mathbf{I})^{-1}
\mathbf{Z}^{\mathsf T}(\mathbf{y}'-\bar{y}'),\qquad y'=2y-1.
\label{eq:ridge}
\end{equation}
Development data were used only for fitting; locked-test data were reserved for final evaluation. Reported quantities were ROC-AUC, balanced accuracy at threshold 0, and class-stratified bootstrap 95\% confidence intervals. Within-domain analyses used 2,000 bootstrap samples by default; cross-generator and duration-sensitivity analyses used 1,000. Univariate comparisons used the Mann--Whitney U test, prompt-paired comparisons used the Wilcoxon signed-rank test, and both applied Benjamini--Hochberg multiple-testing correction.

\section{Complete Operational Record and Preserved Failures}
\label{sec:operations}

\begin{longtable}{p{22mm}p{36mm}p{86mm}}
\caption{Actual operations from design freeze to archival acceptance. Each row is traceable to code, state, logs, or result artifacts.}
\label{tab:timeline}\\
\toprule
Stage & Entry/reference & Operation and verification \\
\midrule
\endfirsthead
\toprule
Stage & Entry/reference & Operation and verification \\
\midrule
\endhead
Design freeze & \outref{O10}, \outref{O11} & Reused the Human/Suno 500+500 cohort and 400/100 split; froze the same Demucs bias; defined the 500+500 extension, 30 s audio standard, and conclusion boundary. \\
Dataset research & \outref{O11}, \coderef{C01} & Compared HeartMuLa-Benchmark, MusicCaps/MusicBench, and Muse. Because the study required vocal-only analysis, Muse metadata was selected for its style fields and time-sectioned lyrics. \\
Prompt construction & \coderef{C01}, \outref{O01} & Stream-filtered Chinese and English records, selected 500 deterministically by salted SHA-256, and independently rebuilt the manifest to verify the same hash. \\
Host selection & \outref{O06} & Independently checked several 5090 hosts and selected 5090-4, whose eight GPUs were idle and whose host load was approximately 0.07. \\
Version freeze & \outref{O07}, \outref{O12} & Frozen Git commits, Hugging Face revisions, pip environments, and weight hashes. Recomputed 87 checkpoint-file hashes after ACE handler initialization. \\
ACE smoke test & \coderef{C02}, \outref{O06} & Used GPU 2. Cold initialization took approximately 32 s and the first track approximately 6.3 s. Output was 30 s, 48 kHz, stereo PCM-16 FLAC. \\
Heart smoke test & \coderef{C03}, \outref{O06} & Used GPUs 0/1. Initialization took 29.52 s and generation 29.43 s. Raw output was 30.08 s, 48 kHz, stereo float WAV. \\
Demucs smoke test & \coderef{C07}, \outref{O05} & Used GPU 3. One standardized ACE item produced two exact 30 s float32 stems. \\
ACE formal generation & \coderef{C02}, \outref{O02} & Processed 500 manifest items, recorded 500 generated events, obtained 500 valid outputs and zero errors. Steady-state generation was approximately 0.95--1.0 s per item. \\
ACE standardization failure & \outref{O04}, \outref{O06} & The first Torchaudio standardization attempt failed for 500/500 files because of a TorchCodec/FFmpeg shared-library mismatch. Failed state and logs were preserved without overwriting successful generation records. \\
ACE standardization recovery & \coderef{C05}, \outref{O04} & Switched to SoundFile plus \texttt{resample\_poly}; 500/500 files succeeded. Minor resampling excursions occurred in 259 songs, totaling 349 samples, and were clipped sample-wise. \\
ACE Demucs & \coderef{C07}, \outref{O05} & Ran 25 batches of 20 tracks, producing 500 separated tracks and 1,000 stems in 1,075.19 s. \\
Initial Heart scheduling & \coderef{C03}, \outref{O03} & Three two-GPU workers processed mutually exclusive index ranges. On detecting a possible boundary race, the workers were interrupted before the target sample was written and restarted with explicit limits. \\
Heart rebalancing & \coderef{C03}, \outref{O03} & After 158 items completed, stopped the old workers and redistributed the actual missing-item list across four two-GPU workers. Valid files were resume-skipped; partial files were removed and retried with the same seeds. \\
Final Heart state & \coderef{C04}, \outref{O03} & Verified 500 files, 500 final valid IDs, zero partial files, and zero final errors. Four BrokenPipe events among 504 historical generated events came from controlled interruption and were recovered. Reconciliation added zero state entries. \\
Heart raw validation & \coderef{C06}, \outref{O04} & A strict $\geq30$ s check failed because 42 native generations ended early; the failure JSON was preserved. A subsequent integrity check allowing 5--30.2 s passed. \\
Heart standardization & \coderef{C05}, \outref{O04} & Zero-padded 42 tracks by 11,501,280 total frames; truncated 457 tracks by at most approximately 80 ms; clipped 5,085 resampling excursions across 164 songs. \\
Heart Demucs & \coderef{C07}, \outref{O05} & Produced 500 separated tracks and 1,000 stems in 594.36 s. \\
Analysis layout & \coderef{C08} & Organized the Human baseline, standardized generated audio, and stems into read-only generator-specific comparison layouts, recording the manifest hash. \\
Primary analysis & \coderef{C09}--\coderef{C12}, \analysisref{A01}--\analysisref{A04} & Computed raw/corrected scalar metrics, frequency vectors, band ablations, univariate effects, and four-panel all/development/locked heatmaps. \\
Paired analysis & \coderef{C13}, \analysisref{A07} & Compared prompt-matched HeartMuLa and ACE-Step tracks using per-pair metric differences, Wilcoxon/BH tests, and average spectral differences. \\
Transfer analysis & \coderef{C15}, \analysisref{A05}, \analysisref{A06} & Ran single-generator train/test matrices, leave-one-generator-out, and pooled-all analyses on both the full cohort and the raw-frozen vocal-active subset. \\
Duration sensitivity & \coderef{C16}, \analysisref{A08} & Excluded 42 short Heart outputs and deterministically sampled an equal number of Human controls within each split before refitting. \\
Archival acceptance & \coderef{C17}, \outref{O08}, \outref{O09} & Checked primary CSVs for NaN/Inf, remote file counts, metadata hashes, core figure layouts, and hashes for 196 lightweight key artifacts. \\
\bottomrule
\end{longtable}

\section{Operational Results and Quality Control}
\label{sec:qc}

\subsection{Generation and Standardization Statistics}

\begin{table}[H]
\centering
\small
\caption{Generation, standardization, and separation results for HeartMuLa and ACE-Step.}
\label{tab:generation-qc}
\begin{tabularx}{\linewidth}{lXX}
\toprule
Item & HeartMuLa & ACE-Step \\
\midrule
Final generation & 500/500; zero effective final errors & 500/500; zero errors \\
Formal generation time & Approximately 69 min 50 s; median 23.34 s/item & Approximately 8 min 9 s; steady state 0.95--1.0 s/item \\
Raw format & 48 kHz stereo float WAV & 48 kHz stereo PCM-16 FLAC \\
Raw duration & Median 30.08 s; 42 below 30 s; minimum 15.12 s & All exactly 30 s \\
Standardized format & 500/500; exact 30 s, 44.1 kHz stereo PCM-16 FLAC & 500/500; same standard \\
Padding/truncation & 42 padded; 457 truncated by at most approximately 80 ms & No duration anomaly \\
Resampling clipping & 164 songs; 5,085 samples & 259 songs; 349 samples \\
Demucs & 500 tracks/1,000 stems; 594.36 s & 500 tracks/1,000 stems; 1,075.19 s \\
\bottomrule
\end{tabularx}
\end{table}

All key counts and format checks are preserved in \outref{O02}--\outref{O05}. The complete remote raw audio, standardized audio, stems, state, and logs are stored under \outref{O09}. The local lightweight archive omits the full audio collection but retains code snapshots, manifests, reports, CSV, PNG, JSON, listening samples, and hash manifests.

\subsection{Vocal-Active Subset}

The vocal-active rule required raw vocal-to-mix RMS of at least $-18$ dB. It was frozen before correction so that correction could not change cohort membership. The development/locked pass counts were Human 249/63, Suno 268/69, HeartMuLa 400/100, and ACE-Step 233/55. The fact that every HeartMuLa item passed is itself a strong domain cue. It may reflect arrangement, vocal proportion, or the response of Demucs to the domain, and must not be interpreted directly as a difference in the underlying generation mechanism.

\section{Primary Spectral Results}
\label{sec:results}

\subsection{Average Spectral-Shape Differences}

\Cref{fig:profiles} summarizes the locked-test, bias-corrected, frame-normalized AI-minus-Human vocal spectral shapes. All three generators are positive in 10--16 kHz, but they diverge strongly in 16--20 kHz.

\begin{figure}[H]
\centering
\includegraphics[width=\linewidth]{../analysis/comparison/locked_corrected_spectral_shape_profiles.png}
\caption{Locked-test mean vocal spectral shapes for the three AI generators relative to Human. A positive value means that AI has greater frame-normalized relative energy at that frequency. Data reference: \texttt{A04}.}
\label{fig:profiles}
\end{figure}

\begin{table}[H]
\centering
\caption{AI-minus-Human band means for the locked-test, bias-corrected, frame-normalized vocal shape.}
\label{tab:band-diff}
\begin{tabular}{lrrrr}
\toprule
Generator & 0.3--5 kHz & 5--10 kHz & 10--16 kHz & 16--20 kHz \\
\midrule
Suno & $-1.158$ & $+1.812$ & $+4.322$ & $+0.146$ \\
HeartMuLa & $-1.469$ & $+1.917$ & $+14.104$ & $-15.912$ \\
ACE-Step & $-2.866$ & $+2.588$ & $+8.465$ & $+10.519$ \\
\bottomrule
\end{tabular}
\end{table}

HeartMuLa forms a steep bandwidth edge near 16.5 kHz. Relative to Human it is approximately 14.1 dB higher in 10--16 kHz but approximately 15.9 dB lower in 16--20 kHz, a pattern consistent with a codec or bandwidth cutoff. ACE-Step is elevated in both 10--16 kHz and 16--20 kHz, more closely resembling persistent upper-band excess. Suno follows a similar direction but with smaller magnitude. The same subjective impression of unnatural high frequencies can therefore arise from opposite physical spectral shapes.

\begin{figure}[H]
\centering
\includegraphics[width=\linewidth]{../analysis/human_vs_heartmula/average_difference_spectrogram_locked_test.png}
\caption{Locked-test average-difference spectrogram for Human and HeartMuLa. The four panels show raw/corrected and absolute/frame-normalized representations. Data reference: \texttt{A02}.}
\label{fig:heart-heatmap}
\end{figure}

\begin{figure}[H]
\centering
\includegraphics[width=\linewidth]{../analysis/human_vs_acestep/average_difference_spectrogram_locked_test.png}
\caption{Locked-test average-difference spectrogram for Human and ACE-Step. ACE-Step's upper-band excess has the opposite direction from the HeartMuLa cutoff above. Data reference: \texttt{A03}.}
\label{fig:ace-heatmap}
\end{figure}

\subsection{Within-Domain Classification}

\begin{table}[H]
\centering
\caption{Best bias-corrected frequency representation on each generator's locked test. Every test set contains 100 Human and 100 AI items.}
\label{tab:single-domain}
\begin{tabular}{llrrr}
\toprule
Generator & Best representation & ROC-AUC & 95\% CI & Balanced accuracy \\
\midrule
Suno & without 5--10 kHz & 0.922 & 0.882--0.957 & 0.855 \\
HeartMuLa & full 0.3--20 kHz & 1.000 & 1.000--1.000 & 1.000 \\
ACE-Step & without 5--10 kHz & 0.988 & 0.971--0.999 & 0.970 \\
\bottomrule
\end{tabular}
\end{table}

\begin{figure}[H]
\centering
\includegraphics[width=0.92\linewidth]{../analysis/comparison/locked_auc_by_generator.png}
\caption{Locked-test band-ablation ROC-AUC for each generator. The 5--10 kHz band alone is discriminative, but the best result generally requires a wider spectral shape. Data reference: \texttt{A04}.}
\label{fig:domain-auc}
\end{figure}

The only-5--10-kHz and only-5--16-kHz AUC values were 0.9999 and 0.9998 for HeartMuLa, 0.9713 and 0.9841 for ACE-Step, and 0.8776 and 0.8992 for Suno. High-frequency information is therefore important, but the best representation in all three domains extends beyond 5--10 kHz. Relative low/mid-band energy, the 10--20 kHz shape, and temporal statistics provide additional information. The complete table is in \analysisref{A04}.

\subsection{Effect of Demucs Correction}

Full-band raw/corrected AUC was 0.879/0.876 for Suno, 1.000/1.000 for HeartMuLa, and 0.9785/0.9792 for ACE-Step. The maximum raw-to-corrected change in a locked average map was approximately 0.260 dB for Suno, 0.274 dB for HeartMuLa, and 0.285 dB for ACE-Step. Thus, the shared average Demucs frequency response is not the principal source of discrimination. This result \textbf{does not} rule out different separator errors for Human audio and for individual generators.

\subsection{Sensitivity to Short HeartMuLa Outputs}

Forty-two of 500 raw HeartMuLa generations were shorter than 30 s; zero-padding could therefore become a shortcut. \coderef{C16} excluded these items and deterministically sampled an equal number of Human examples within the development and locked splits, yielding 366+366 development and 92+92 locked items. Full-band corrected AUC and balanced accuracy both remained 1.000. Trailing silence is therefore insufficient to explain the saturated result, although this test does not remove genre, mastering, or separator-domain confounding. Detailed CSV and metadata are in \analysisref{A08}.

\section{Cross-Generator Generalization and Generator Fingerprints}
\label{sec:cross}

\subsection{Single-Source Train/Test Matrix}

If a representation encoded a universal AI mechanism, it would preserve score direction and ranking across generators. The observed transfer did not satisfy this requirement. For the corrected representation without 5--10 kHz on the full cohort, Suno$\rightarrow$HeartMuLa achieved 0.905, whereas HeartMuLa$\rightarrow$Suno achieved 0.379. HeartMuLa$\rightarrow$ACE-Step achieved 0.326, whereas ACE-Step$\rightarrow$HeartMuLa achieved 0.593. AUC below 0.5 indicates that the score direction reverses after domain shift, rather than merely becoming less accurate.

\begin{figure}[H]
\centering
\includegraphics[width=\linewidth]{../analysis/cross_generator_vocal_active/cross_generator_auc_heatmaps.png}
\caption{Cross-generator Human-versus-AI ROC-AUC on the raw-frozen vocal-active subset. Near-perfect diagonal entries and degraded or nearly reversed off-diagonal entries show that high within-domain AUC includes generator-specific fingerprints. Data reference: \texttt{A06}.}
\label{fig:cross-heatmaps}
\end{figure}

\subsection{Leave-One-Generator-Out and Pooled Training}

\begin{table}[H]
\centering
\caption{Best corrected locked-test AUC under leave-one-generator-out evaluation. The target generator is absent from training.}
\label{tab:logo}
\begin{tabular}{lrrl}
\toprule
Target generator & Full tracks & Vocal-active & Vocal-active 95\% CI \\
\midrule
Suno & 0.698 & 0.702 & 0.612--0.785 \\
HeartMuLa & 0.949 & 0.942 & 0.891--0.978 \\
ACE-Step & 0.834 & 0.935 & 0.882--0.980 \\
\bottomrule
\end{tabular}
\end{table}

With each locked Human item included only once, pooled-all training achieved AUC 0.943 (0.920--0.963) on the full cohort and 0.969 (0.950--0.987) on the vocal-active cohort. The vocal-active pooled model achieved 0.940, 0.998, and 0.988 on Suno, HeartMuLa, and ACE-Step, respectively. These results support a transferable broadband/high-frequency heuristic. However, leave-one-out transfer to unseen Suno is much weaker than transfer to the other two systems; the current system should not be described as a universal detector.

\section{Prompt-Paired HeartMuLa--ACE-Step Analysis}
\label{sec:paired}

The two models share the same 500 style descriptions, lyrics, languages, splits, and integer seeds. Their direct comparison therefore controls text conditions more closely than either comparison with Human. The stochastic processes are not equivalent, however, and Demucs may still respond differently to their generated timbres.

\begin{table}[H]
\centering
\caption{Prompt-paired HeartMuLa and ACE-Step features on the bias-corrected locked test. H-A is the median within-pair HeartMuLa-minus-ACE-Step difference.}
\label{tab:paired}
\begin{tabular}{lrrrr}
\toprule
Metric & HeartMuLa & ACE-Step & H-A & Separation AUC \\
\midrule
Sibilance burst rate & 1.625 & 1.130 & $+0.527$ & 0.904 \\
HF crest (dB) & 14.07 & 16.82 & $-2.826$ & 0.857 \\
HF modulation share & 0.2465 & 0.1417 & $+0.0943$ & 0.770 \\
Fakeprint periodicity & 0.0327 & 0.0531 & $-0.0286$ & 0.757 \\
HF tilt 5--16 kHz & $-6.13$ & $-8.06$ & $+3.083$ & 0.707 \\
HF ratio 5--16 kHz & $-3.639$ & $-6.305$ & $+3.184$ & 0.695 \\
Air ratio 12--20 kHz & $-13.85$ & $-19.80$ & $+5.947$ & 0.686 \\
\bottomrule
\end{tabular}
\end{table}

\begin{figure}[H]
\centering
\includegraphics[width=\linewidth]{../analysis/paired_generators/average_difference_spectrogram_locked_test.png}
\caption{Prompt-paired locked-test average vocal difference for HeartMuLa minus ACE-Step. The shared correction approximately cancels in the direct dB comparison; the maximum change in the locked absolute map is only approximately 0.271 dB. Data reference: \texttt{A07}.}
\label{fig:paired-map}
\end{figure}

The strongest paired differences are sibilance bursts, HF crest, and high-frequency envelope modulation, rather than a single static energy measurement. This supports an applied filter that combines frequency distribution with temporal texture instead of a fixed equalization threshold.

\section{MIDI-Grounded Validation of a Music-Understanding Model}
\label{sec:music-flamingo-probe}

\subsection{Motivation and Admission Criterion}

The proposed applied filter includes two non-spectral hypotheses: AI-generated
arrangements may exhibit reduced dynamics variation in MIDI-controlled
instruments, and they may exhibit less varied rhythmic development. These
hypotheses can be measured deterministically when aligned MIDI is available.
For waveform-only deployment, a music-understanding model might appear to offer
a convenient semantic measurement. This section therefore tests Music Flamingo
\citep{musicflamingo2025} before allowing any model-derived judgment to enter
the Human--AI detector.

This is a capability-validation experiment, not an AI-versus-Human
classification test. The model feature is admitted only if its lower 95\%
Wilson bound exceeds 0.5, physical-excerpt selection is at least 0.75 consistent
under A/B reversal, exact JSON parse rate is at least 0.95, performance remains
above chance among low-nuisance pairs, and no direction reversal appears across
major styles. Failure leaves the deterministic MIDI/MIR measurement in the
heuristic set and excludes the language-model judgment.

The official Music Flamingo evaluations cover broad captioning, music QA,
instrument, lyric, and reasoning tasks, but do not report dedicated accuracy
for MIDI velocity-variation or bar-to-bar rhythm-variation judgments
\citep{musicflamingohf2026}. MF-Skills was not used for evaluation because its
card describes synthetic/model-assisted annotations and it was used in Music
Flamingo training \citep{mfskills2025}.

\subsection{Open Dataset and Objective Targets}

The open Groove MIDI Dataset (GMD) supplies human-performed drum MIDI, aligned
synthesized audio, velocity and timing, tempo/meter metadata, and an official
test split under CC BY 4.0 \citep{gmd2019}. Only 4/4 \texttt{beat}
performances in the test split were eligible. The official MIDI-only archive
defined every target. Audio was transported through a 16-kHz Hugging Face
mirror at frozen revision
\texttt{a505eb614310d60f6844042709c286b8a4d1988f}; mirror labels were not used
\citep{gmdmirror2026}.

For each pitch $p$ with at least four note onsets, dynamics variation combines
two onset-weighted measurements:
\begin{equation}
  D_p^{(1)} = \frac{\sigma(v_p)}{\mu(v_p)}, \qquad
  D_p^{(2)} = \frac{\operatorname{median}_i |v_{p,i}-v_{p,i-1}|}
  {\operatorname{median}_i v_{p,i}}.
\end{equation}
The primary dynamics score is the equal-weight mean of the within-corpus robust
$z$-scores of the two pitch-conditional summaries. Pitch conditioning prevents
changes in the mixture of kick, snare, cymbal, and tom events from masquerading
as expressive dynamics.

For rhythm, onsets are quantized to a 16th-note grid and mapped to coarse drum
classes. Each of four bars is represented by a binary class-by-grid vector
$b_j$. The rhythm-variation score is the average pairwise Jaccard distance:
\begin{equation}
  R = \frac{1}{\binom{4}{2}}\sum_{i<j}
  \left(1-\frac{|b_i\cap b_j|}{|b_i\cup b_j|}\right).
\end{equation}

The final pool contained 116 eligible four-bar windows from 36 official test
performances. Strict within-performance matching yielded 22 dynamics pairs and
10 rhythm pairs. Dynamics target gaps were 0.531--4.148 (median 0.990), with a
median rhythm-nuisance gap of 0.044 and a median RMS gap of 0.282 dB. Rhythm
target gaps were 0.085--0.255 (median 0.109), with a median dynamics-nuisance
gap of 0.175 and a median RMS gap of 0.244 dB. The shortfall from the intended
32 rhythm pairs was retained rather than weakening matching rules after data
inspection.

\subsection{Pairwise Protocol and 5090 Inference}

Every pair came from the same source performance. Because the released chat
template emits one \texttt{<sound>} placeholder per user turn, excerpt A and
excerpt B were concatenated with one second of silence. Paired durations were
14.913--25.000 s (median 19.824 s), below the 30-s processor window. Each pair
was queried in both orders; the correct physical excerpt is unchanged while its
letter flips.

Only the seed-randomized first ordering of each pair enters primary accuracy and
its Wilson interval, avoiding dependence between an item and its reversal.
Reversed queries measure physical-choice order consistency and positional bias.
\texttt{TIE}, invalid JSON, and runtime failures count as incorrect.

Inference used \path{nvidia/music-flamingo-2601-hf} revision
\MFHash{6b5be086d52f65a1e204cb0faf70bf54e2741ecd}
\citep{musicflamingocheckpoint2026}, bfloat16, greedy decoding, and at most 64
new tokens. It ran only on physical GPU 1 of host \texttt{5090-3}, an NVIDIA
GeForce RTX 5090 exposed as process-local \texttt{cuda:0}. The runtime was
PyTorch 2.11.0+cu128 and Transformers 5.16.1. The 16,534,531,504-byte weights
used LFS SHA-256/OID
\MFHash{4da26452d3978503fb064023a69f7867a1f235a2a5a89d3d9146ebf2e04eb15b}.
Peak allocated memory was 15.62 GiB. The 64 primary queries required 17.53 s
after load, with median generation latency 0.216 s.

\subsection{Results}

\begin{table}[H]
\centering
\caption{Preregistered Music Flamingo capability results. Accuracy uses one
seed-randomized ordering per physical pair; response diagnostics use both
orders. TIE is incorrect for accuracy.}
\label{tab:mf-primary}
\resizebox{\linewidth}{!}{%
\begin{tabular}{lrrrrr}
\toprule
Task & Pairs & Accuracy (95\% Wilson CI) & TIE rate & JSON parse & Order consistency \\
\midrule
Dynamics & 22 & 0.000 (0.000--0.149) & 1.000 & 1.000 & 0.000 \\
Rhythm   & 10 & 0.300 (0.108--0.603) & 0.500 & 1.000 & 0.000 \\
\bottomrule
\end{tabular}
}
\end{table}

All 44 dynamics queries returned \texttt{TIE}. Among 20 rhythm queries, ten
returned \texttt{TIE} and ten returned A; none returned B. Accuracy in the
smaller-nuisance half was 0.000 for dynamics and 0.400 for rhythm. Rhythm
correctness did not increase with target separation ($\rho=-0.114$,
$p=0.754$). Both features fail the complete admission rule. Order consistency
is defined as selecting the same physical excerpt in both orders; a TIE/TIE
pair is a stable abstention but not a valid ranking.

\begin{figure}[H]
  \centering
  \includegraphics[width=0.97\linewidth]{\MFFigureRoot/analysis_comparison/primary_vs_forced_choice.png}
  \caption{Primary and post-hoc forced-choice results. The left panel can be
  misleading in isolation: the apparent forced-choice dynamics accuracy is
  generated by a universal first-position response. The right panel exposes
  the TIE collapse and the A-only positional shortcut.}
  \label{fig:mf-primary-forced}
\end{figure}

\subsection{Post-hoc Forced-Choice Sensitivity}

After the primary TIE collapse was observed, one explicitly post-hoc prompt
removed TIE and required A or B; no audio, order, target, checkpoint, or decode
setting changed. Music Flamingo selected A on all 64 queries. This produced an
apparent dynamics accuracy of 15/22 = 0.682 (95\% CI 0.473--0.836) and rhythm
accuracy of 3/10 = 0.300 (0.108--0.603). Both had a first-position selection
rate of 1.000 and physical-choice order consistency of 0.000. The 0.682 value
is therefore the proportion of randomized first orders whose correct excerpt
happened to occupy position A, not evidence of dynamics perception. This
post-hoc analysis cannot replace the preregistered result.

\subsection{Decision for the Applied Filter}

Music Flamingo is excluded from the current dynamics and rhythm feature set. It
obeyed the JSON output format perfectly, but did not produce a usable,
order-invariant perceptual ranking. The thesis retains the pitch-conditional
velocity statistics and bar-pattern Jaccard distance as deterministic metrics
when MIDI is available. For waveform-only deployment, the appropriate next
features are onset-conditioned energy dispersion, adjacent-onset energy change,
bar-synchronous pattern novelty, onset-density variation, tempogram change, and
self-similarity novelty on percussion or accompaniment stems. These MIR proxies
must be calibrated against aligned MIDI before entering the generator-disjoint
Human--AI ablation.

The result does not show that all music-understanding models are incapable of
these judgments, nor does it measure AI-versus-Human accuracy. GMD is
percussion-only and uses synthesized electronic-drum audio; only ten rhythm
pairs passed strict matching. A future preregistered replication should use a
new corpus, balance A/B labels exactly, freeze a forced-choice prompt in
advance, add calibrated human judgments, and include aligned pitched-instrument
audio/MIDI such as piano before generalizing to bass, guitar, voice, or
orchestral parts.

\subsection{Reproducible References}

\begin{longtable}{p{42mm}p{104mm}}
\caption{Music Flamingo probe artifact index. Paths are relative to the
experiment's lightweight artifact root.}\\
\toprule
Artifact & Purpose \\
\midrule
\endfirsthead
\toprule
Artifact & Purpose \\
\midrule
\endhead
\path{EXPERIMENT_PROTOCOL.md} & Frozen questions, metrics, pairing, admission rule, and two pre-output implementation repairs. \\
\path{DATASET_MODEL_AUDIT.md} & Model/dataset coverage, leakage analysis, and source selection. \\
\path{code/prepare_gmd_probe.py} & MIDI parsing, objective measurements, matching, and paired-audio materialization. \\
\path{code/run_music_flamingo_probe.py} & Resume-safe 5090-only model inference and raw output capture. \\
\path{code/analyze_probe.py} & Accuracy, Wilson intervals, order checks, style tables, and primary figure. \\
\path{code/compare_conditions.py} & Primary-versus-forced comparison and response-distribution figure. \\
\path{pairs.csv}; \path{prompt_manifest.jsonl} & Frozen physical pairs and 64 order-specific queries. \\
\path{state/inference.jsonl} & All preregistered responses, labels, parse status, correctness, and timing. \\
\path{state/inference_forced_choice.jsonl} & All post-hoc forced-choice responses. \\
\path{analysis/}; \path{analysis_forced_choice/} & CSV/JSON summaries, style cells, and figures for both conditions. \\
\path{analysis_comparison/condition_comparison.csv} & Direct numerical source for \cref{fig:mf-primary-forced}. \\
\path{logs/} & Formal logs, package freeze, and GPU inventory. \\
\bottomrule
\end{longtable}

The complete audio-bearing archive is retained at
\path{/mnt/nfs-code/users/yi/music_flamingo_midi_rhythm_probe_20260903};
the local thesis workspace contains the lightweight manifests, raw model
responses, code, statistics, logs, and figures.


\section{External Validation of Dynamics, Rhythm, and Structure Heuristics}
\label{sec:external-mir-benchmark}

The Music Flamingo result in the preceding section rejects an audio-language-model shortcut, but it does not establish that deterministic MIR alternatives are valid. A separate benchmark was therefore frozen before consulting the AI-versus-Human development data. The complete protocol, deployment addendum, provenance audit, commands, and recovery log are indexed as \outref{O14}--\outref{O17}. This experiment tests \emph{measurement validity}, not AI detection: a passing result would mean that a feature measures its claimed musical property on an external annotated corpus, not that the feature separates Human and AI music.

\subsection{Datasets, models, and admission policy}

Table~\ref{tab:external-mir-design} separates the three base tasks from the detector-facing summaries derived from their outputs. Dataset selection used salted deterministic ordering and was fixed before inference. Learned inference was restricted to GPU 0 of an NVIDIA GeForce RTX 5090 on \texttt{5090-5}; deterministic preprocessing, scoring, and 10,000-replicate recording/track-cluster bootstrap ran on the same host.

\begin{table}[H]
\centering
\caption{Independent feature-validation design. ``External'' is relative to the released model or to this thesis detector; no AI/Human labels are used in these benchmarks.}
\label{tab:external-mir-design}
\begin{tabularx}{\linewidth}{p{24mm}p{32mm}X p{28mm}}
\toprule
Family & External data & Base estimator and frozen sample & Detector-facing gate \\
\midrule
Rhythm/bar & GTZAN & Beat This! \texttt{final0/1/2}; 993 scored tracks per seed; native non-DBN decoding \citep{beatthisrepo2026} & $\rho\geq0.75$, CI lower $\geq0.65$ \\
Dynamics & MAESTRO v3 test & 50 composition-disjoint recordings; 200 frozen 30 s windows; aligned MIDI velocity \citep{maestrodataset2026} & $\rho\geq0.50$, CI lower $>0.35$ \\
Section/bar length & SALAMI & 44 tracks, 68 macro-annotation files; All-In-One boundaries plus Beat This! downbeats \citep{salamirepo2026,allinonerepo2026} & boundary F1@3 s $\geq0.60$, then $\rho\geq0.70$ \\
\bottomrule
\end{tabularx}
\end{table}

The GTZAN archive had 999 annotation files; six jazz files contained only a beat column and were explicitly excluded by the released downbeat dataloader. MAESTRO transport used a fixed Hugging Face mirror revision, but all 50 WAV and 50 MIDI files were verified against official Google ZIP member names, uncompressed sizes, and CRC-32 values, with 100/100 accepted. SALAMI selection probed all 476 official Internet Archive mappings, found 49 reachable records, and froze all 44 records that also met duration and annotation requirements. The manifests and transport evidence are in \outref{O15}.

The original All-In-One implementation depends on a legacy NATTEN interface unavailable on the RTX 5090. The run therefore used \texttt{openmirlab/all-in-one-infer} commit \texttt{3c93b4a}, which preserves the released architecture/configuration and eight author checkpoints while providing a pure-PyTorch neighborhood-attention backend. This compatibility deviation was recorded before batch inference and did not change data, weights, thresholds, or scoring.

\subsection{Rhythm localization succeeds, rhythm richness does not}

Beat This! passes both timing gates. Across \texttt{final0/1/2}, mean beat F1 is 0.891 (SD 0.004) and mean downbeat F1 is 0.783 (SD 0.005). For \texttt{final0}, track-bootstrap beat F1 is 0.892 (95\% CI 0.879--0.904) and downbeat F1 is 0.787 (0.768--0.806). Beat CMLt/AMLt is 0.798/0.903 and downbeat CMLt/AMLt is 0.678/0.799.

However, none of the five second-order summaries reaches the detector gate (Table~\ref{tab:rhythm-feature-fidelity}). This is not a contradiction: an event tracker can localize most beats while small timing errors, missing events, and downbeat-number errors strongly perturb variation and entropy estimates.

\begin{table}[H]
\centering
\caption{GTZAN reference--prediction fidelity of detector-facing rhythm/bar summaries using Beat This! \texttt{final0}. All confidence intervals use 10,000 track-bootstrap replicates.}
\label{tab:rhythm-feature-fidelity}
\begin{tabular}{lrrl}
\toprule
Summary & Spearman $\rho$ & 95\% CI & Decision \\
\midrule
Robust IBI CV & 0.377 & 0.321--0.430 & Reject \\
Tempo total variation & 0.224 & 0.161--0.285 & Reject \\
Tempo entropy & 0.532 & 0.482--0.580 & Reject \\
Bar off-mode fraction & 0.145 & 0.064--0.213 & Reject \\
Bar-length entropy & 0.144 & 0.067--0.210 & Reject \\
\bottomrule
\end{tabular}
\end{table}

Five tracks were excluded from rhythm-pair fidelity because the prediction contained fewer than 16 beats. Bar-summary fidelity excluded 139 tracks: 133 references and 61 predictions contained fewer than eight complete bars, with overlap between reasons. These counts are preserved rather than imputed. Direct outputs are indexed as \analysisref{A10}.

\subsection{MIDI dynamics are recoverable only with oracle support}

The primary MAESTRO experiment evaluates three non-learned acoustic estimators at aligned MIDI note onsets and pitches. Broadband and harmonic onset contrasts fail. In contrast, post-onset harmonic attack level passes: median within-$(\text{recording},\text{pitch})$ $\rho$ is 0.815 (95\% CI 0.800--0.834), and low-versus-high velocity quartile ranking accuracy is 0.948 (0.941--0.956). With oracle onset and pitch, all three 30 s summaries also pass (Table~\ref{tab:dynamics-oracle-deployment}).

A deployed detector does not receive MIDI onset and pitch. Before running the deployment analysis, a separate addendum therefore froze a waveform-only spectral-flux onset detector and harmonic-peak attack estimator. It obtains onset F1 0.751, matched-event $\rho=0.626$ (0.594--0.659), and quartile-ranking accuracy 0.905 (0.890--0.919). Thus, individual attack strength is measurable. Nevertheless, no 30 s dynamics-variety summary reaches the pre-specified window-level gate.

\begin{table}[H]
\centering
\caption{MAESTRO window-level dynamics fidelity. Oracle values use aligned MIDI onset/pitch to select the acoustic measurement; waveform-only values use no MIDI input. Constant +6 dB gain drift is below $3.6\times10^{-15}$ for every row.}
\label{tab:dynamics-oracle-deployment}
\begin{tabular}{lrrrl}
\toprule
Feature & Oracle $\rho$ & Waveform $\rho$ & Waveform 95\% CI & Detector decision \\
\midrule
90--10\% dynamics span & 0.857 & 0.414 & 0.250--0.561 & Reject \\
Dynamics IQR & 0.826 & 0.358 & 0.204--0.495 & Reject \\
Adjacent attack change & 0.798 & $-0.067$ & $-0.221$--0.091 & Reject \\
\bottomrule
\end{tabular}
\end{table}

This distinction prevents an oracle-to-deployment validity error. The supported claim is narrow: piano attack strength can be estimated from the waveform. The unsupported claim is that the present aggregation measures dynamics variety over a polyphonic 30 s passage faithfully enough for the applied filter. Oracle and waveform outputs are indexed as \analysisref{A11}.

\subsection{Section and bar-length variety fail external validation}

All-In-One plus Beat This! produced complete outputs for 44/44 selected SALAMI tracks. Across 68 annotator rows, macro-boundary F1 is 0.342 (0.289--0.395) at 0.5 s and 0.511 (0.464--0.557) at 3 s. The latter fails both the 0.60 point threshold and 0.55 lower-bound threshold. The derived section-duration and bar-length summaries also have low fidelity: duration median $\rho=0.042$, duration CV $-0.017$, duration entropy 0.065, bars median 0.258, bars CV 0.194, and bar off-mode fraction 0.327. All six are rejected. Functional section labels were not used. Direct outputs are indexed as \analysisref{A12}.

\begin{figure}[H]
\centering
\includegraphics[width=\linewidth]{../../dynamics_rhythm_external_benchmark_20260903/figures/external_benchmark_dashboard.png}
\caption{External validity dashboard. Error bars are 95\% recording/track-cluster bootstrap confidence intervals. Panels A and the event-level dynamics result show that a base task may pass even when the downstream heuristic in panels B--D fails. Artifact reference: \texttt{A13}.}
\label{fig:external-mir-dashboard}
\end{figure}

\subsection{Gate decision for the applied filter}

\noindent\fcolorbox{black!25}{codegray}{%
\parbox{0.96\linewidth}{\textbf{Decision: zero waveform-deployable features are admitted.} Consequently, no AI/Human detector ablation was run for these three heuristic families, the Human/Suno/HeartMuLa/ACE-Step development data were not used for rescue selection, and the already examined external AI/Human test was not touched. The machine-readable decision is \analysisref{A13}.}}

This negative result is methodologically useful. Running an AI/Human ablation after observing failed external fidelity would turn the detector dataset into a tuning set for a measurement that has not shown construct validity. The most promising next hypothesis concerns dynamics aggregation rather than event estimation: instrument-aware event clustering, non-vocal sub-stems, and bar-synchronous robust summaries. Such changes require a new frozen external benchmark before detector application.


\section{Controlled-Change Validation and Detector Extension}
\label{sec:change-extension}

\subsection{Why a Second External Test Was Required}

The annotation-fidelity benchmark in the preceding section asks whether an
audio estimator reproduces independently annotated dynamics, tempo variation,
or section-length statistics.  The subsequent experimental request poses a
related but different operational question: can the complete waveform
workflow detect a deliberately introduced change on external music, and does
the identically defined output then add information to the existing AI/Human
classifier?  The controlled-change protocol in \outref{O18} was frozen before
audio transformation or detector-cohort extraction.  The original fidelity
results remain unchanged; the present test appends rather than replaces them.

Three paired interventions were created by \coderef{C26}.  Fifty MAESTRO test
recordings contributed an original 30 s piano excerpt, a constant $-6$ dB
negative control, and a smoothly varying $\pm6$ dB version.  Forty-four SALAMI
recordings contributed an original excerpt, a constant 0.8x-speed negative
control, and a version alternating 0.8x and 1.25x speed every four source
seconds.  The same SALAMI recordings also contributed two four-excerpt
mosaics: one with four equal 7.5 s units and one with 3, 5, 8, and 14 s units.
All variants were stereo, 44.1 kHz, and exactly 30 s.  Source recordings were
not modified.

The waveform dynamics extractor measures the robust span, interquartile
range, and adjacent change of onset-centered harmonic attack strength.  Beat
This! \texttt{final0} supplies beat times for IBI CV, normalized local-tempo
total variation, and tempo entropy.  All-In-One \texttt{harmonix-all} supplies
section boundaries, while Beat This! downbeats convert section duration into
bar-count variety.  The family statistic is the equal-weight average of
robustly scaled paired feature differences.  The frozen sensitivity gate
requires probability of superiority at least 0.75 and a 10,000-replicate
paired-bootstrap 95\% lower bound greater than 0.60.  Code and raw outputs are
indexed as \coderef{C27}--\coderef{C31} and \analysisref{A14}.

\subsection{External Controlled-Change Results}

\begin{table}[H]
\centering
\caption{External controlled-change sensitivity.  PASS/FAIL uses the frozen
family-level gate; individual feature results do not override it.}
\label{tab:controlled-change}
\begin{tabular}{lrrrrl}
\toprule
Family & Pairs & Median composite change & $P(\Delta>0)$ & 95\% CI & Gate \\
\midrule
Dynamics & 50 & 0.619 & 0.880 & [0.780, 0.960] & PASS \\
Rhythm change & 44 & 1.226 & 0.705 & [0.580, 0.818] & FAIL \\
Section/bar variety & 44 & 0.000 & 0.523 & [0.466, 0.591] & FAIL \\
\bottomrule
\end{tabular}
\end{table}

Dynamics is therefore the only externally sensitivity-qualified family.  Its
span and IQR features had superiority probabilities 0.960 and 0.880,
respectively, whereas the constant-gain composite shifted in the opposite
direction ($P=0.340$).  The rhythm composite did not pass, although tempo
entropy alone was strongly responsive ($P=0.917$, 95\% CI
[0.806, 1.000]); this is a candidate for a new confirmatory corpus, not a
post-hoc admission in the present experiment.  The constant-speed rhythm
composite was materially smaller and negative ($P=0.318$).

The structure failure was primarily an observability failure.  Only one of 44
paired mosaics yielded a complete paired set of detector-facing section/bar
statistics.  All-In-One commonly represented a 30 s mosaic with one long
same-label segment plus edge fragments, so the four intended lengths were not
available to the summary.  Missing pair contributions were neutral in the
predeclared composite; they did not create a favorable result.  The failed
threshold and low coverage are retained explicitly rather than repaired using
AI/Human labels.

\subsection{Application to the Frozen 2,000-Track Cohort}

The exact workflows were next applied to 500 Human FMA, 500 Suno, 500
HeartMuLa, and 500 ACE-Step tracks with the existing 400/100
development/locked split.  Beat This! produced 1,998 files; two Human tracks
had no output and 14 further Human files were empty.  All-In-One produced
2,000 JSON files.  Every learned inference process ran on RTX 5090 GPUs;
deterministic aggregation ran on CPU.  A symlink-only cohort manifest avoided
rewriting any standardized audio.

The dynamics signal needs one important alignment with the motivating
MIDI-controlled-instrument hypothesis.  The external MAESTRO signal contains
no vocal.  Before inspecting any detector result, the detector-side signal was
therefore frozen as the sample-wise sum of Demucs \texttt{bass},
\texttt{drums}, and \texttt{other}, excluding \texttt{vocals}.  All 2,000
tracks had complete non-vocal stems and dynamics measurements.  This prevents
the earlier vocal high-frequency cues from simply being counted again as an
instrument-dynamics result, although Demucs domain response remains a possible
confound.

A fixed ridge linear classifier ($\alpha=10$) used development-only medians,
standardization, and missingness indicators.  The same locked examples were
used for every paired AUC comparison.  The directly recomputed 24-feature
bias-corrected vocal spectral representation is the baseline.  The only
sensitivity-qualified extension is \texttt{spectral+dynamics}; rhythm,
structure, and \texttt{all-new} results are reported as exploratory because
their external family gates failed.  Evaluation code, per-track scores, and
bootstrap tables are \coderef{C32} and \analysisref{A15}.

\begin{table}[H]
\centering
\footnotesize
\caption{Locked ROC-AUC of the externally qualified dynamics extension and
the exploratory all-new representation.  Delta and CI compare
\texttt{spectral+dynamics} with the spectral baseline on identical tracks.}
\label{tab:change-detector-extension}
\resizebox{\linewidth}{!}{%
\begin{tabular}{lrrrrr}
\toprule
Scenario & Dynamics only & Spectral & Spectral+dyn. & Paired $\Delta$ [95\% CI] & Spectral+all-new \\
\midrule
Within Suno & 0.701 & 0.787 & 0.807 & $+0.020$ [0.005, 0.036] & 0.857 \\
Within HeartMuLa & 0.917 & 0.999 & 1.000 & $+0.000$ [$-0.001$, 0.002] & 0.999 \\
Within ACE-Step & 0.850 & 0.950 & 0.954 & $+0.003$ [$-0.004$, 0.011] & 0.970 \\
Pooled all & 0.816 & 0.875 & 0.892 & $+0.017$ [0.006, 0.027] & 0.914 \\
Leave out Suno & 0.672 & 0.659 & 0.707 & $+0.047$ [0.023, 0.076] & 0.747 \\
Leave out HeartMuLa & 0.899 & 0.944 & 0.962 & $+0.018$ [0.005, 0.034] & 0.981 \\
Leave out ACE-Step & 0.853 & 0.657 & 0.721 & $+0.064$ [0.043, 0.087] & 0.741 \\
\bottomrule
\end{tabular}
}
\end{table}

The qualified result is positive but modest in the pooled evaluation:
spectral+dynamics increases AUC from 0.8754 to 0.8919, a paired difference of
0.0165 (95\% CI [0.0059, 0.0272]).  More importantly, the increment is positive
in every generator-held-out transfer: 0.0472 for Suno, 0.0179 for HeartMuLa,
and 0.0641 for ACE-Step, with all three intervals above zero.  This pattern is
consistent with additional cross-generator information rather than only a
single-generator fingerprint.  Dynamics alone is also informative, reaching
pooled AUC 0.8156 and held-out AUC 0.8528 for ACE-Step.  Within HeartMuLa and
ACE-Step increments are small because the spectral baselines are already near
their within-domain ceilings.

The exploratory union reaches pooled AUC 0.9139 and improves every held-out
scenario, but it is not an admitted detector layer: rhythm and structure failed
the external family gate.  Structure alone is weak (pooled AUC 0.5663), and its
pooled spectral increment is 0.0049 with a 95\% interval crossing zero.  Its
complete four-feature coverage is also source-dependent: only 7\% of locked
Human, 12\% of locked Suno, 7\% of locked HeartMuLa, and 4\% of locked ACE-Step
tracks.  The classifier can therefore learn availability artifacts through
the missingness indicators.  Rhythm coverage is much higher (90\%--100\% on
locked domains), and its exploratory pooled spectral increment is 0.0175, but
the failed external family gate prevents a confirmatory authorship claim.

\begin{figure}[H]
\centering
\includegraphics[width=\linewidth]{../../dynamics_rhythm_change_detector_extension_20260903/figures/change_detector_extension_dashboard.png}
\caption{Two-stage detector extension.  Panel A shows the external
controlled-change family gates; panel B shows pooled locked AUC; panel C shows
the exploratory all-new AUC increment.  Only dynamics is externally
sensitivity-qualified.  Data references: \texttt{A14--A16}.}
\label{fig:change-detector-extension}
\end{figure}

\subsection{Interpretation for the Applied Filter}

The deployment recommendation changes in one narrow way.  Non-vocal dynamics
may be enabled as a separately reported, low-weight extension to the vocal
spectral layer because it passes an external change-sensitivity gate and adds
locked cross-generator information.  The system should expose its three
contributing measurements and a valid-stem flag, not present them as direct
evidence of MIDI authorship.  Rhythm entropy is promising but remains a
candidate feature pending a new independently frozen rhythm benchmark.
Section/bar-length variety remains disabled for decision making on 30 s
windows because the estimator did not observe the controlled changes reliably
and its coverage is strongly domain-dependent.

These results still do not establish a universal AI-music detector.  The
external controlled transformations validate sensitivity to the specified
interventions, not construct specificity; a compressor, human quantization,
genre, arrangement, mastering, or separator behavior can produce similar
values.  The existing independent AI/Human external test was not reopened for
feature selection.  The new extension should therefore be frozen and tested
there, or preferably on a new generator-disjoint and production-matched test,
only after this thesis-stage analysis is finalized.


\section{Frozen External Test Set: 500 Tracks per Non-Development Generator}
\label{sec:external-generator-500}

To test whether the applied filter transfers beyond Suno, HeartMuLa, and
ACE-Step, a new evaluation-only set was frozen from AIME
\citep{aime2025,aimehf2026}. AIME was selected because it provides Human and
generated music under the same 500 three-tag prompt conditions and exposes
exactly 500 outputs for every included generator. The source was pinned to main
revision
\hashvalue{b84d4be5eda830b6eb714998569dba73530f2601} and parquet revision
\hashvalue{1bdacac93127439e361bdd19d575d8b596bca4e3}. Neither detector scores nor
audio-derived features were used to choose membership.

\subsection{Generator and Leakage Boundaries}

The resulting manifest contains 5,000 AI tracks from ten generators, exactly
500 per generator, and 500 shared MTG-Jamendo Human controls. Suno v3, Suno
v3.5, generic Suno labels, HeartMuLa, and ACE-Step are hard exclusions. The
primary cohort comprises six generator families absent from the earlier
100-track AIME diagnostic; the secondary cohort expands four families that
appeared in that diagnostic. The exact 50 AI and 50 Human tracks previously
used are marked row by row and are not treated as pristine holdout examples.

\begin{table}[H]
\centering
\small
\caption{Frozen external AI cohorts. Every listed model contributes exactly 500 tracks.}
\label{tab:external-generator-cohorts}
\begin{tabularx}{\linewidth}{lXrr}
\toprule
Cohort & Generators & Models & AI tracks \\
\midrule
Primary pristine generator-disjoint & MusicGen Small, MusicGen Medium, MusicGen Large, AudioLDM 2 Large, AudioLDM 2 Music, Mustango & 6 & 3,000 \\
Secondary family-seen expansion & Udio, Riffusion, Stable Audio v1, Stable Audio v2 & 4 & 2,000 \\
Shared external Human control & MTG-Jamendo & 1 & 500 \\
\bottomrule
\end{tabularx}
\end{table}

All 500 prompt conditions were assigned with seed \texttt{20260904} to five
disjoint groups of 100. For every generator, a dedicated view contains its 500
AI tracks and the same 500 Human tracks; each of its five group manifests
contains 100 AI and 100 Human tracks with identical condition-ID sets. This
supports per-generator reporting and prompt-condition-paired bootstrap
intervals without selecting favorable random mixtures after evaluation.

\subsection{Standardisation and Native-Band Guard}

Every selected recording was converted to a deterministic 10.000 s crop at
44.1 kHz, stereo, signed 16-bit FLAC by \coderef{C33}. No zero-padding,
looping, loudness normalization, peak normalization, equalization, denoising,
or source separation was applied. Ten seconds is the longest no-padding
duration shared by all selected AIME systems; consequently, this set is
appropriate for frame-level spectral tests but not for long-horizon structural
claims.

Resampling does not create valid high-frequency evidence. The manifest therefore
records original sample rate, native Nyquist frequency, and per-band eligibility
for every track. The detector must mask features above the original Nyquist
limit rather than interpreting upsampled silence as an AI cue.

\begin{table}[H]
\centering
\caption{Native bandwidth of the 5,500-track external set before standardisation.}
\label{tab:external-native-bandwidth}
\begin{tabular}{lrr}
\toprule
Quantity & Tracks & Fraction \\
\midrule
Original sample rate 16 kHz & 1,500 & 27.3\% \\
Original sample rate 32 kHz & 1,500 & 27.3\% \\
Original sample rate 44.1 kHz & 500 & 9.1\% \\
Original sample rate 48 kHz & 2,000 & 36.4\% \\
Eligible through 8 kHz & 5,500 & 100.0\% \\
Eligible through 10 kHz & 4,000 & 72.7\% \\
Eligible through 20 kHz & 2,500 & 45.5\% \\
\bottomrule
\end{tabular}
\end{table}

\subsection{Integrity and Confirmatory Evaluation Policy}

The builder froze membership before materialisation. The frozen-selection and
final-manifest SHA-256 digests are, respectively,
\begin{center}
\small
\hashvalue{5f25ed368276784b0eb5cc8d56c97e381e4c6c3875eec71f6f29cd6cc9b5363d}\\
\hashvalue{49dbe01b8be98857361b3a437db44ee937698e014bf70ccc7d443c24ad342fb4}.
\end{center}

All 5,500 standardised files passed exhaustive decoding, duration, format,
path, and SHA-256 verification; there were 5,500 unique IDs, 5,500 unique output
hashes, no duplicate raw-audio hash groups, and zero verification errors. The
source audit \coderef{C34} covered 202 retained parquet shards totaling
53,386,356,734 bytes. Although retained boundary shards contained 500 Suno v3
and 260 Suno v3.5 rows, the allowlist rejected every one of them. The complete
protocol, source audit, manifests, checksum ledger, and execution logs are
indexed by \outref{O21}--\outref{O23}.

No detector score has yet been computed on this collection. The already frozen
detector and thresholds must first be evaluated on the six-model primary cohort,
with per-generator and pooled metrics reported. The four family-seen models must
then be reported separately, both including and excluding exact prior-diagnostic
overlaps. Any representation or threshold changed after seeing these results is
a new hypothesis and requires another untouched test set.



\section{Integrated Conclusions and Applied-Filter Design}
\label{sec:conclusion}

\subsection{Conclusions Supported by the Current Evidence}

\begin{enumerate}
  \item \textbf{The 10--16 kHz region is the most consistent anomalous band across the three current AI domains.} Each generator has a positive frame-normalized mean relative to Human, though the magnitude differs.
  \item \textbf{Anomalies in 16--20 kHz must be detected in both directions.} HeartMuLa has a sharp cutoff above approximately 16.5 kHz, whereas ACE-Step has an upper-band excess. A rule that equates more high-frequency energy with AI would miss HeartMuLa.
  \item \textbf{Static spectral shape and dynamic texture should be combined.} HF tilt, band ratios, crest, modulation, sibilance bursts, HNR, and pitch confidence describe complementary phenomena.
  \item \textbf{The frozen average Demucs correction does not change the main conclusions.} This weakens the explanation that a shared average separator response creates the result, but it does not exclude class-dependent error.
  \item \textbf{Within-domain AUC substantially overstates universality.} Cross-generator score reversal demonstrates sensitivity to codec, bandwidth, and generation-pipeline fingerprints.
  \item \textbf{Music Flamingo is not a validated dynamics or rhythm feature.} The MIDI-grounded capability probe produced 0/22 dynamics accuracy, 3/10 rhythm accuracy, and zero physical-choice order consistency.
  \item \textbf{External fidelity and controlled-change sensitivity answer different questions.} The original detector-facing summaries fail the strict annotation-fidelity gate, but the subsequently frozen end-to-end manipulation test qualifies the non-vocal dynamics family. Its addition raises pooled spectral AUC from 0.8754 to 0.8919 and produces positive, non-zero-crossing increments for all three held-out generators. Rhythm and section/bar-length family gates still fail, so their favorable cohort ablations remain exploratory.
\end{enumerate}

\subsection{Recommended Layered Applied Filter}

Given the current evidence, the thesis prototype should be organized into four layers instead of treating a within-domain ridge classifier as the final detector:

\begin{enumerate}
  \item \textbf{Cross-model layer:} compute the 0.3--20 kHz frame-normalized shape, 5--16 kHz tilt and ratio, 10--16 kHz excess, sibilance bursts, HNR, and pitch confidence on the Demucs vocal stem. Train on multiple generators in development, then evaluate on locked and unseen generators.
  \item \textbf{Fingerprint layer:} explicitly detect bandwidth cutoffs, 16--20 kHz excess, comb-like periodicity, HF crest, and HF modulation. This layer should answer which known generator a sample resembles, rather than directly equating resemblance with AI status.
  \item \textbf{Qualified non-vocal dynamics layer:} sum the Demucs bass, drums, and other stems; report attack-level span, IQR, and adjacent change with a validity flag. Use this as a low-weight extension because its pooled increment is modest and separator/production confounding remains possible.
  \item \textbf{Deferred rhythm and structure layer:} retain tempo entropy as a candidate for a new confirmatory corpus, but do not activate the current rhythm family. Disable All-In-One section/bar-length variety for 30 s decisions because the external gate and complete-feature coverage fail. Music Flamingo judgments remain disabled.
\end{enumerate}

The final system should report a continuous score, the principal contributing bands, a vocal-activity flag, and a domain/out-of-distribution warning. A hard real-world decision threshold is not justified until a matched Human cohort and unseen-generator validation are available.

\subsection{Required Limitations and Next Experiments}

The 500 Muse prompts are strongly biased toward Pop (319), Ballad (292), Female Vocal (286), Piano (263), and C-pop (233). The 500 Human FMA items are dominated by Rock (163), Electronic (115), Experimental (45), and Hip-Hop (42), with only 21 Pop items. Current AUC can therefore contain genuine generator fingerprints together with genre, language, vocal-arrangement, mastering, and provenance differences. The next priorities are:

\begin{enumerate}
  \item construct a Human cohort matched as closely as possible in style, language, vocal activity, and mastering;
  \item evaluate the frozen six-generator pristine AIME cohort first, without changing the current representation or threshold, and report all six per-generator results rather than only a pooled score;
  \item test robustness under MP3/AAC compression, resampling, equalization, loudness normalization, and mastering chains;
  \item repeat the key comparisons using oracle vocals or multiple separators to estimate class-dependent separation error; and
  \item align the spectral measurements with blinded listening labels to test whether they actually explain the perceived high-frequency cue;
  \item replicate the language-model capability test on a new, exactly A/B-balanced corpus with a forced-choice prompt frozen in advance;
  \item repeat the qualified non-vocal dynamics result with oracle or alternative stems, compression/mastering perturbations, and a new production-matched generator-disjoint test;
  \item freeze tempo entropy alone on a new external corpus before considering its favorable exploratory detector increment; and
  \item replace or lengthen the section estimator before re-testing section/bar variety, since 30 s complete-feature coverage is only 4--12\% in most locked domains.
\end{enumerate}

\section{Reproduction and Integrity}
\label{sec:repro}

Exact reproduction commands are indexed as \outref{O12}. The complete remote experiment root is
\begin{center}
\filepath{/mnt/nfs-code/users/yi/open_models_spectral_500_20260901}
\end{center}
and the existing Human/Suno/Demucs-bias baseline is
\begin{center}
\filepath{/mnt/nfs-code/users/yi/demucs_bias_corrected_1000_20260901}.
\end{center}

The local lightweight archive root is
\begin{center}
\filepath{/Users/yi/Documents/code/music/artifacts/open_models_spectral_500_20260901}.
\end{center}

The complete Music Flamingo capability experiment and its lightweight thesis archive are, respectively,
\begin{center}
\filepath{/mnt/nfs-code/users/yi/music_flamingo_midi_rhythm_probe_20260903}\\
\filepath{/Users/yi/Documents/code/music/artifacts/music_flamingo_midi_rhythm_probe_20260903}.
\end{center}

The complete external MIR feature benchmark and its local archive are, respectively,
\begin{center}
\filepath{/mnt/nfs-code/users/yi/dynamics_rhythm_external_benchmark_20260903}\\
\filepath{/Users/yi/Documents/code/music/artifacts/dynamics_rhythm_external_benchmark_20260903}.
\end{center}

The complete controlled-change and detector-extension experiment and its local
lightweight archive are, respectively,
\begin{center}
\filepath{/mnt/nfs-code/users/yi/dynamics_rhythm_change_detector_extension_20260903}\\
\filepath{/Users/yi/Documents/code/music/artifacts/dynamics_rhythm_change_detector_extension_20260903}.
\end{center}

The complete external 500-per-generator test set and its local lightweight
archive are, respectively,
\begin{center}
\filepath{/mnt/nfs-code/users/yi/external_generator_500_testset_20260904}\\
\filepath{/Users/yi/Documents/code/music/artifacts/external_generator_500_testset_20260904}.
\end{center}

The lightweight manifest \outref{O08} covers 196 key files. Its own SHA-256 digest is:
\begin{center}
\small\hashvalue{abe0b6b6705aec5f00cd8d516e207a4209428df9e833356d2f13ab2712b8c272}
\end{center}

\clearpage
\appendix
\setlength{\tabcolsep}{5pt}

\section{Code Reference Index}
\label{app:code}

Unless otherwise stated, paths in the following table are relative to the local lightweight archive root. The main text links to these records through identifiers such as \coderef{C01}. Archived code snapshots keep the reported results tied to the executed implementation even if the working copy later changes.

\begin{longtable}{p{15mm}p{55mm}p{80mm}}
\caption{Code artifact index.}\\
\toprule
ID & Relative path & Function \\
\midrule
\endfirsthead
\toprule
ID & Relative path & Function \\
\midrule
\endhead
\codeanchor{C01} & \filepath{code/experiment/build_muse_manifest.py} & Filters Muse metadata, constructs lyrics for the first 30 s, deterministically selects 500 items, and assigns splits and seeds. \\
\codeanchor{C02} & \filepath{code/experiment/run_acestep_batch.py} & Initializes ACE-Step, performs frozen-parameter inference, writes atomically, resumes safely, and records per-item state and hashes. \\
\codeanchor{C03} & \filepath{code/experiment/run_heartmula_batch.py} & Runs HeartMuLa/HeartCodec on separate GPUs with ranged workers, resume logic, state, and hashes. \\
\codeanchor{C04} & \filepath{code/experiment/reconcile_generated_state.py} & Read-only verification of existing audio after controlled interruption and reconstruction of missing state entries without changing audio. \\
\codeanchor{C05} & \filepath{code/experiment/standardize_generated_audio.py} & SoundFile decoding, polyphase resampling, stereo/30 s/PCM-16 FLAC standardization, and clip/pad/truncate accounting. \\
\codeanchor{C06} & \filepath{code/experiment/validate_generated_batch.py} & Validates manifest/state/audio integrity, sample rate, channels, duration, and finite samples. \\
\codeanchor{C07} & \filepath{code/experiment/run_demucs_generated.py} & RTX 5090 Demucs batch wrapper, frozen htdemucs parameters, stem metadata verification, and state logging. \\
\codeanchor{C08} & \filepath{code/experiment/prepare_analysis_layout.py} & Builds analysis layouts and manifests for the Human baseline and each generator's audio and stems. \\
\codeanchor{C09} & \filepath{code/analysis_lib/stem_spectral_ablation.py} & Implements the STFT, 22 interpretable scalar metrics, active-frame rules, statistics, and ridge model. \\
\codeanchor{C10} & \filepath{code/analysis_lib/analyze_demucs_artifacts.py} & Computes Demucs artifact metrics, frequency vectors, band ablations, bootstrap intervals, and locked evaluation. \\
\codeanchor{C11} & \filepath{code/analysis_lib/analyze_bias_corrected_vocals.py} & Applies the frozen Demucs bias filter and compares raw/corrected metrics, univariate effects, and classifications. \\
\codeanchor{C12} & \filepath{code/experiment/analyze_average_class_spectrogram.py} & Produces 128-band population-average difference spectrograms and four-panel heatmaps. \\
\codeanchor{C13} & \filepath{code/experiment/analyze_paired_generators.py} & Computes prompt-paired HeartMuLa/ACE metrics, Wilcoxon/BH results, and difference heatmaps. \\
\codeanchor{C14} & \filepath{code/experiment/compare_generators.py} & Builds cross-generator summary tables, AUC figures, metric effects, and spectral profiles. \\
\codeanchor{C15} & \filepath{code/experiment/cross_generator_generalization.py} & Runs single-source matrices, leave-one-generator-out, pooled-all, and active-only analyses. \\
\codeanchor{C16} & \filepath{code/experiment/heartmula_duration_robust_ablation.py} & Excludes short HeartMuLa outputs and deterministically samples equal-size Human controls for sensitivity analysis. \\
\codeanchor{C17} & \filepath{code/experiment/collect_run_metadata.py} & Collects host, Git, environment, and key-file hashes into run metadata. \\
\codeanchor{C18} & \filepath{../dynamics_rhythm_external_benchmark_20260903/code/prepare_maestro_manifest.py} & Deterministically selects 50 MAESTRO test recordings and freezes four 30 s windows per recording. \\
\codeanchor{C19} & \filepath{../dynamics_rhythm_external_benchmark_20260903/code/verify_maestro_transport.py} & Verifies all transported MAESTRO WAV/MIDI files against official ZIP member size and CRC-32 metadata. \\
\codeanchor{C20} & \filepath{../dynamics_rhythm_external_benchmark_20260903/code/analyze_rhythm_benchmark.py}; \filepath{../dynamics_rhythm_external_benchmark_20260903/code/aggregate_rhythm_results.py} & Scores Beat This! event metrics, detector-facing richness fidelity, exclusions, bootstrap intervals, and three-seed gates. \\
\codeanchor{C21} & \filepath{../dynamics_rhythm_external_benchmark_20260903/code/analyze_maestro_dynamics.py} & Evaluates three acoustic dynamics estimators with oracle MIDI onsets/pitches and gain-invariant window summaries. \\
\codeanchor{C22} & \filepath{../dynamics_rhythm_external_benchmark_20260903/code/analyze_maestro_waveform_dynamics.py} & Implements the deployment-addendum waveform-only onset, harmonic attack, matching, and clustered-bootstrap analysis. \\
\codeanchor{C23} & \filepath{../dynamics_rhythm_external_benchmark_20260903/code/prepare_salami_manifest.py} & Audits official SALAMI Internet Archive URLs and freezes the reachable eligible sample independently of model output. \\
\codeanchor{C24} & \filepath{../dynamics_rhythm_external_benchmark_20260903/code/analyze_salami_structure.py} & Scores All-In-One boundaries and second/bar length fidelity across all annotator files. \\
\codeanchor{C25} & \filepath{../dynamics_rhythm_external_benchmark_20260903/code/build_external_benchmark_report.py} & Aggregates all external gates and generates the machine-readable decision, report, and dashboard. \\
\codeanchor{C26} & \filepath{../dynamics_rhythm_change_detector_extension_20260903/code/prepare_external_change_audio.py} & Creates deterministic MAESTRO gain, SALAMI tempo, and equal/variable section-length controlled audio. \\
\codeanchor{C27} & \filepath{../dynamics_rhythm_change_detector_extension_20260903/code/workflow_features.py} & Implements shared waveform-dynamics, beat-derived rhythm, and boundary/downbeat structure features. \\
\codeanchor{C28} & \filepath{../dynamics_rhythm_change_detector_extension_20260903/code/analyze_external_change.py} & Runs paired robust composites, 10,000-replicate bootstrap intervals, negative controls, and frozen family gates. \\
\codeanchor{C29} & \filepath{../dynamics_rhythm_change_detector_extension_20260903/code/build_detector_manifest.py}; \filepath{../dynamics_rhythm_change_detector_extension_20260903/code/make_inference_shards.py} & Builds the exact 2,000-row symlink cohort and verifies disjoint balanced RTX 5090 inference shards. \\
\codeanchor{C30} & \filepath{../dynamics_rhythm_change_detector_extension_20260903/code/extract_detector_features.py} & Extracts non-vocal dynamics and joins frozen beat, downbeat, and section outputs without dropping missing rows. \\
\codeanchor{C31} & \filepath{../dynamics_rhythm_change_detector_extension_20260903/code/build_report.py}; \filepath{../dynamics_rhythm_change_detector_extension_20260903/code/build_integrity_manifest.py} & Builds the bilingual evidence reports, dashboard, and lightweight SHA-256 manifest. \\
\codeanchor{C32} & \filepath{../dynamics_rhythm_change_detector_extension_20260903/code/evaluate_detector_extension.py} & Fits the fixed ridge representation, evaluates within/pooled/held-out scenarios, and bootstraps same-track AUC increments. \\
\codeanchor{C33} & \filepath{../external_generator_500_testset_20260904/code/build_aime_500_testset.py} & Freezes allowed AIME rows, imports prior-overlap flags, assigns prompt groups, materialises deterministic clips, constructs paired model views, and exhaustively verifies all outputs. \\
\codeanchor{C34} & \filepath{../external_generator_500_testset_20260904/code/audit_aime_source.py} & Hashes the retained revision-pinned parquet shards and records shard-level model counts and byte totals. \\
\bottomrule
\end{longtable}

\section{Output and Process Reference Index}
\label{app:outputs}

Except for \outref{O09}, the paths below are relative to the local lightweight archive root.

\begin{longtable}{p{15mm}p{58mm}p{77mm}}
\caption{Generation, standardization, separation, and audit-output index.}\\
\toprule
ID & Path & Contents \\
\midrule
\endfirsthead
\toprule
ID & Path & Contents \\
\midrule
\endhead
\outanchor{O01} & \filepath{prompts/prompt_manifest.jsonl}; \filepath{prompts/prompt_manifest_metadata.json}; \filepath{prompts/selected_source_rows.jsonl} & Frozen 500-item prompt set with split, language, seed, source fields, and hashes. \\
\outanchor{O02} & \filepath{state/acestep_generation.jsonl}; \filepath{logs/acestep_generation.log} & Per-item ACE-Step status, duration, format, byte count, SHA-256, and primary log. \\
\outanchor{O03} & \filepath{state/heartmula_generation*.jsonl}; \filepath{logs/heartmula_generation*.log} & Evidence for HeartMuLa multi-stage workers, controlled interruption, resume/rebalance, and final status. \\
\outanchor{O04} & \filepath{reports/acestep_standardization_summary.json}; \filepath{reports/heartmula_standardization_summary.json}; \filepath{reports/*generation_validation*.json}; \filepath{state/acestep_standardization_failed_torchcodec.jsonl} & Standardization counts, raw validation, short Heart outputs, and the failed ACE TorchCodec attempt. \\
\outanchor{O05} & \filepath{state/demucs_acestep.jsonl}; \filepath{state/demucs_heartmula.jsonl}; \filepath{logs/demucs_*.log} & Demucs batches, elapsed time, cumulative validation counts, and smoke/formal logs. \\
\outanchor{O06} & \filepath{RUN_LOG.md}; \filepath{logs/*pip_freeze.txt}; \filepath{logs/*smoke.log} & Host selection, environment, smoke tests, scheduling, failures, performance, and final status. \\
\outanchor{O07} & \filepath{logs/heartmula_weight_sha256.txt}; \filepath{logs/acestep_weight_sha256_final.txt} & Frozen weight hashes and per-file hashes for final ACE checkpoints after handler initialization. \\
\outanchor{O08} & \filepath{reports/LIGHTWEIGHT_SHA256.txt}; \filepath{reports/RUN_METADATA.json} & SHA-256 for 196 lightweight key artifacts and run metadata. \\
\outanchor{O09} & \filepath{/mnt/nfs-code/users/yi/open_models_spectral_500_20260901} & Complete remote raw/standardized audio, 2,000 stems, state, logs, models, and analysis. \\
\outanchor{O10} & \filepath{EXPERIMENT_PROTOCOL.md} & Pre-frozen questions, parameters, operational rules, and conclusion boundaries. \\
\outanchor{O11} & \filepath{MODEL_DATASET_RESEARCH.md}; \filepath{THESIS_RESULTS_CN.md} & Model/dataset selection evidence and the earlier Chinese-stage conclusions. \\
\outanchor{O12} & \filepath{REPRODUCE_COMMANDS.md} & Exact generation, standardization, validation, Demucs, Human-versus-AI, paired, and cross-generator commands. \\
\outanchor{O13} & \filepath{samples/README.md}; \filepath{samples/*} & HeartMuLa/ACE mixes and vocal stems for two locked prompts, retained for listening review. \\
\outanchor{O14} & \filepath{../dynamics_rhythm_external_benchmark_20260903/EXPERIMENT_PROTOCOL.md}; \filepath{../dynamics_rhythm_external_benchmark_20260903/DEPLOYMENT_VALIDATION_ADDENDUM.md}; \filepath{../dynamics_rhythm_external_benchmark_20260903/DATASET_MODEL_AUDIT.md} & Frozen external hypotheses, thresholds, waveform-only addendum, source revisions, licenses, exclusions, and RTX 5090 compatibility deviation. \\
\outanchor{O15} & \filepath{../dynamics_rhythm_external_benchmark_20260903/manifests/} & Frozen MAESTRO and SALAMI selections, URL audit, exact windows, and 100/100 official-ZIP transport verification. \\
\outanchor{O16} & \filepath{../dynamics_rhythm_external_benchmark_20260903/RUN_LOG.md}; \filepath{../dynamics_rhythm_external_benchmark_20260903/REPRODUCE_COMMANDS.md}; \filepath{../dynamics_rhythm_external_benchmark_20260903/logs/} & Commands, environment, GPU restriction, failures, resume history, inference logs, and integrity evidence. \\
\outanchor{O17} & \filepath{/mnt/nfs-code/users/yi/dynamics_rhythm_external_benchmark_20260903} & Complete remote source checkouts, audio, checkpoints, environment, intermediate predictions, logs, and final results. \\
\outanchor{O18} & \filepath{../dynamics_rhythm_change_detector_extension_20260903/EXPERIMENT_PROTOCOL.md} & Frozen change-sensitivity design plus timestamped outcome-handling and detector-side non-vocal signal addenda. \\
\outanchor{O19} & \filepath{../dynamics_rhythm_change_detector_extension_20260903/RUN_LOG.md}; \filepath{../dynamics_rhythm_change_detector_extension_20260903/REPRODUCE_COMMANDS.md}; \filepath{../dynamics_rhythm_change_detector_extension_20260903/logs/} & Complete execution history, recovery/re-sharding record, software versions, commands, and inference logs. \\
\outanchor{O20} & \filepath{/mnt/nfs-code/users/yi/dynamics_rhythm_change_detector_extension_20260903} & Complete remote controlled audio, 2,000 four-stem directories, model intermediates, predictions, results, and reports. \\
\outanchor{O21} & \filepath{../external_generator_500_testset_20260904/EXPERIMENT_PROTOCOL.md}; \filepath{../external_generator_500_testset_20260904/DATASET_AUDIT.md}; \filepath{../external_generator_500_testset_20260904/TESTSET_REPORT.md}; \filepath{../external_generator_500_testset_20260904/RUN_LOG.md}; \filepath{../external_generator_500_testset_20260904/REPRODUCE_COMMANDS.md} & Frozen evaluation rules, source-selection rationale, cohort interpretation, executed run history, and exact reproduction commands. \\
\outanchor{O22} & \filepath{../external_generator_500_testset_20260904/testset/}; \filepath{../external_generator_500_testset_20260904/source/audit/}; \filepath{../external_generator_500_testset_20260904/logs/} & Frozen/final manifests, paired per-model and group views, native-band QC, source-shard checksum ledger, summaries, and complete execution logs. \\
\outanchor{O23} & \filepath{/mnt/nfs-code/users/yi/external_generator_500_testset_20260904} & Complete remote 53.4 GB source mirror, 3.5 GB standardised audio, code, state, manifests, audits, and logs. \\
\bottomrule
\end{longtable}

\section{Data-Analysis Result Reference Index}
\label{app:analysis}

The following paths are relative to the local lightweight archive root. Each directory contains a human-readable \texttt{REPORT.md}; CSV/JSON/NPZ files are numerical sources, while PNG files are derived visualizations.

\begin{longtable}{p{15mm}p{60mm}p{75mm}}
\caption{Data-analysis artifact index.}\\
\toprule
ID & Path & Associated result or data \\
\midrule
\endfirsthead
\toprule
ID & Path & Associated result or data \\
\midrule
\endhead
\analysisanchor{A01} & \filepath{analysis/human_vs_suno_average/} & Human-versus-Suno all/development/locked average-difference maps and metadata. \\
\analysisanchor{A02} & \filepath{analysis/human_vs_heartmula/} & HeartMuLa raw/corrected metrics, band classifications, univariate effects, vocal activity, and average maps. \\
\analysisanchor{A03} & \filepath{analysis/human_vs_acestep/} & Complete corresponding raw/corrected ACE-Step results. \\
\analysisanchor{A04} & \filepath{analysis/comparison/} & Three-generator classification comparison, band summary, metric effects, AUC figures, and spectral profiles. \\
\analysisanchor{A05} & \filepath{analysis/cross_generator/} & Full-cohort single-source matrix, leave-one-out, pooled-all CSV/reports, and heatmaps. \\
\analysisanchor{A06} & \filepath{analysis/cross_generator_vocal_active/} & Cross-generator results for the raw-frozen vocal-active sensitivity cohort. \\
\analysisanchor{A07} & \filepath{analysis/paired_generators/} & Metrics for 500 HeartMuLa-versus-ACE pairs, Wilcoxon/BH tests, and four-panel difference maps. \\
\analysisanchor{A08} & \filepath{analysis/heartmula_duration_robust/} & Duration sensitivity after excluding 42 short outputs, using 366+366/92+92 items. \\
\analysisanchor{A09} & \filepath{analysis/demucs_frequency_bias.npz} & Demucs frequency bias pre-frozen on external MUSDB oracle vocals. \\
\analysisanchor{A10} & \filepath{../dynamics_rhythm_external_benchmark_20260903/results/rhythm_benchmark/} & Three-seed Beat This! metrics, final0 bootstrap intervals, exclusions, feature fidelity, and zero rhythm/bar admissions. \\
\analysisanchor{A11} & \filepath{../dynamics_rhythm_external_benchmark_20260903/results/maestro_dynamics/}; \filepath{../dynamics_rhythm_external_benchmark_20260903/results/maestro_waveform_dynamics/} & Oracle and waveform-only MIDI-dynamics validation, event measurements, window features, bootstrap intervals, and gain tests. \\
\analysisanchor{A12} & \filepath{../dynamics_rhythm_external_benchmark_20260903/results/salami_structure/} & Boundary metrics and six section/bar-length fidelity results over 44 tracks and 68 annotation files. \\
\analysisanchor{A13} & \filepath{../dynamics_rhythm_external_benchmark_20260903/results/aggregate_summary.json}; \filepath{../dynamics_rhythm_external_benchmark_20260903/results/admission_decisions.csv}; \filepath{../dynamics_rhythm_external_benchmark_20260903/figures/external_benchmark_dashboard.png}; \filepath{../dynamics_rhythm_external_benchmark_20260903/reports/REPORT_EN.md} & Aggregate zero-admission decision, dashboard, and English evidence report. \\
\analysisanchor{A14} & \filepath{../dynamics_rhythm_change_detector_extension_20260903/external_results/controlled_change_features.csv}; \filepath{../dynamics_rhythm_change_detector_extension_20260903/external_results/controlled_change_sensitivity.csv}; \filepath{../dynamics_rhythm_change_detector_extension_20260903/external_results/controlled_change_summary.json} & Per-variant features, individual and family paired effects, bootstrap intervals, and sensitivity gates. \\
\analysisanchor{A15} & \filepath{../dynamics_rhythm_change_detector_extension_20260903/detector_results/evaluation/detector_metrics.csv}; \filepath{../dynamics_rhythm_change_detector_extension_20260903/detector_results/evaluation/incremental_auc_deltas.csv}; \filepath{../dynamics_rhythm_change_detector_extension_20260903/detector_results/evaluation/detector_scores.csv} & Locked within-generator, pooled, and held-out metrics plus paired spectral-extension AUC changes and per-track scores. \\
\analysisanchor{A16} & \filepath{../dynamics_rhythm_change_detector_extension_20260903/detector_results/new_features.csv}; \filepath{../dynamics_rhythm_change_detector_extension_20260903/detector_results/evaluation/feature_coverage.csv}; \filepath{../dynamics_rhythm_change_detector_extension_20260903/figures/change_detector_extension_dashboard.png}; \filepath{../dynamics_rhythm_change_detector_extension_20260903/reports/REPORT_EN.md} & Joined feature coverage, complete evidence report, and controlled-change/detector dashboard. \\
\bottomrule
\end{longtable}

\section{Mapping Reported Numbers to Original Artifacts}
\label{app:claim-map}

\begingroup
\footnotesize
\renewcommand{\arraystretch}{1.05}
\begin{longtable}{p{57mm}p{35mm}p{58mm}}
\caption{Direct data sources for key claims.}\\
\toprule
Claim & Reference & Direct file \\
\midrule
\endfirsthead
\toprule
Claim & Reference & Direct file \\
\midrule
\endhead
Three-generator bandwise AI-minus-Human dB & \analysisref{A04} & \filepath{analysis/comparison/average_spectrogram_band_summary.csv} \\
Within-domain AUC/CI/balanced accuracy & \analysisref{A04} & \filepath{analysis/comparison/generator_classification_comparison.csv} \\
Strongest scalar cues for each model & \analysisref{A02}, \analysisref{A03} & \filepath{univariate_effects.csv} and \filepath{REPORT.md} \\
Leave-one-out and pooled AUC & \analysisref{A05}, \analysisref{A06} & \filepath{cross_generator_evaluation.csv} \\
HeartMuLa/ACE paired metrics & \analysisref{A07} & \filepath{paired_metric_summary.csv}; \filepath{paired_vocal_metrics.csv} \\
Sensitivity to 42 short HeartMuLa outputs & \analysisref{A08}, \outref{O04} & \filepath{classification_ablation.csv}; \filepath{heartmula_standardization_summary.json} \\
Maximum raw/corrected map change & \analysisref{A01}--\analysisref{A03} & \filepath{average_difference_metadata.json} \\
GPU, timing, failure, and recovery history & \outref{O02}--\outref{O06} & JSONL state, smoke/formal logs, and \filepath{RUN_LOG.md} \\
Version, weight, and archive integrity & \outref{O07}, \outref{O08} & Weight hashes, \filepath{LIGHTWEIGHT_SHA256.txt}, and \filepath{RUN_METADATA.json} \\
Beat/downbeat timing versus rhythm-richness fidelity & \analysisref{A10} & \filepath{rhythm_benchmark_summary.json}; per-seed CSV and source summaries \\
Oracle versus waveform-only dynamics validity & \analysisref{A11} & \filepath{summary.json}; \filepath{window_feature_fidelity.csv}; \filepath{window_features.csv} \\
SALAMI boundary and section/bar-length fidelity & \analysisref{A12} & \filepath{summary.json}; \filepath{feature_fidelity.csv}; \filepath{per_annotator_metrics.csv} \\
External gate and no-ablation decision & \analysisref{A13}, \outref{O14}--\outref{O16} & \filepath{aggregate_summary.json}; \filepath{admission_decisions.csv}; protocols and run log \\
Controlled-change sensitivity and negative controls & \analysisref{A14}, \outref{O18} & \filepath{controlled_change_sensitivity.csv}; \filepath{controlled_change_summary.json} \\
Qualified dynamics and exploratory rhythm/structure AUC increments & \analysisref{A15} & \filepath{detector_metrics.csv}; \filepath{incremental_auc_deltas.csv}; \filepath{detector_scores.csv} \\
Feature missingness and 30 s section observability & \analysisref{A16} & \filepath{feature_coverage.csv}; \filepath{new_features.csv}; dashboard and report \\
External 500-per-generator membership, bandwidth, and integrity & \outref{O21}--\outref{O23} & \filepath{frozen_selection_summary.json}; \filepath{materialization_summary.json}; \filepath{source_audit_summary.json}; manifests, views, and verification log \\
\bottomrule
\end{longtable}
\endgroup

\clearpage
\bibliographystyle{unsrtnat}
\bibliography{references}

\end{document}
